\documentclass[../../main.tex]{subfiles}% 注意这里的文件路径不能用 ./main.tex ,否则用latexmk编译子文件会报错
\graphicspath{{\subfix{./image/}}{\subfix{../../image/}}} % 指定图片引用路径,后续可以直接使用图片文件名
% 如果图片存放在上级目录,则使用 ../../image/ 作为备用路径

% 例如:
% \begin{figure}[H]
% \centering
% \includegraphics[width=0.3\textwidth]{图.png}
% \caption{}
% \label{figure:图}
% \end{figure}
% 注意:上述\label{}一定要放在\caption{}之后,否则引用图片序号会只会显示??.

\begin{document}

\section{正则性结果的证明}\label{section:ch07-正则证明}

\begin{proof}[\cref{proposition:ch07-7.1.2}的证明]
(i) 只证\eqref{align:ch07-7.15}，\eqref{align:ch07-7.16}类似。设 $\widetilde\zeta\in C^1([0,T];C^1_c(\mathbb{R}))$，$\widetilde\zeta(\cdot,T)=0$ 且 $\widetilde\zeta\geqslant0$。在\eqref{align:ch07-7.9}中令 $l\to\infty$，易得
\begin{align}
\int_0^T\langle u_s(\cdot,t),\partial_t\widetilde\zeta(\cdot,t)\rangle\,\mathrm{d}t\geqslant-\langle u_{0s},\widetilde\zeta(\cdot,0)\rangle. \label{align:ch07-7.157}
\end{align}
由 $C([0,T];C_c(\mathbb{R}))$ 的可分性，存在零测集 $N\subseteq(0,T)$，使对所有 $\zeta\in C([0,T];C_c(\mathbb{R}))$ 与 $\tau\in(0,T)\setminus N$，有
\begin{align}
\lim_{h\to0^+}\frac1h\int_\tau^{\tau+h}\langle u_s(\cdot,t),\zeta(\cdot,t)\rangle\,\mathrm{d}t
=\langle u_s(\cdot,\tau),\zeta(\cdot,\tau)\rangle. \label{align:ch07-7.158}
\end{align}
设 $t_1,t_2\in(0,T)\setminus N$，$0<t_1<t_2<T$。由标准逼近论证，可在\eqref{align:ch07-7.157}中取 $\widetilde\zeta(x,t)=g_h(t)\zeta(x,t)$，其中
\begin{align*}
g_h(t)\triangleq \frac1h(t-t_1)\chi_{[t_1,t_1+h]}(t)+\chi_{(t_1+h,t_2)}(t)+\frac1h(t_2+h-t)\chi_{[t_2,t_2+h]}(t),
\end{align*}
$h\in(0,\min\{t_2-t_1,T-t_2\})$。令 $h\to0^+$ 并利用\eqref{align:ch07-7.158}，得
\begin{align}
\langle u_s(\cdot,t_2),\zeta(\cdot,t_2)\rangle
\leqslant\int_{t_1}^{t_2}\langle u_s(\cdot,t),\partial_t\zeta(\cdot,t)\rangle\,\mathrm{d}t
+\langle u_s(\cdot,t_1),\zeta(\cdot,t_1)\rangle. \label{align:ch07-7.159}
\end{align}
在\eqref{align:ch07-7.159}中取 $\zeta(x,t)=\rho(x)$（$\rho\in C_c(\mathbb{R})$，$\rho\geqslant0$），即得\eqref{align:ch07-7.15}。

(ii) 如 (i) 那样论证，存在零测集 $N\subseteq(0,T)$，使对所有 $\rho\in C_c(\mathbb{R})$ 与 $\tau\in(0,T)\setminus N$，有
\begin{align}
\lim_{h\to0^+}\frac1h\int_\tau^{\tau+h}\langle u(\cdot,t),\rho\rangle\,\mathrm{d}t=\langle u(\cdot,\tau),\rho\rangle. \label{align:ch07-7.160}
\end{align}
置 $\Omega_j\triangleq [-j-1,-j]\cup[j,j+1]$，$Q_{j,\tau}\triangleq \Omega_j\times(0,\tau)$（$\tau\in(0,T)$，$j\in\mathbb{N}$）。设 $\{\rho_j\}\subseteq C^1_c(\mathbb{R})$ 满足 $\rho_j=1$ 在 $[-j,j]$，$\operatorname{supp}\rho_j\subseteq[-j-1,j+1]$，$0\leqslant\rho_j\leqslant1$ 且 $|\rho_j'|\leqslant2$。在\eqref{align:ch07-7.7}中取 $\zeta(x,t)=f_h(t)\rho_j(x)$，其中
\begin{align*}
f_h(t)\triangleq \chi_{[0,\tau)}(t)+\frac1h(\tau+h-t)\chi_{[\tau,\tau+h]}(t),
\end{align*}
$h$ 充分小。则
\begin{align*}
\iint_{Q_{j,\tau}}f_h(t)\phi(u_r)(x,t)\rho_j'(x)\,\mathrm{d}x\,\mathrm{d}t
=\frac1h\int_\tau^{\tau+h}\langle u(\cdot,t),\rho_j\rangle\,\mathrm{d}t-\langle u_0,\rho_j\rangle.
\end{align*}
令 $h\to0^+$ 并利用\eqref{align:ch07-7.160}，得
\begin{align}
\iint_{Q_{j,\tau}}\phi(u_r)(x,t)\rho_j'(x)\,\mathrm{d}x\,\mathrm{d}t
=\langle u(\cdot,\tau),\rho_j\rangle-\langle u_0,\rho_j\rangle. \label{align:ch07-7.161}
\end{align}
因 $|\phi(u_r)|\leqslant Mu_r$（见\eqref{align:ch07-7.1}）且 $u_r\in L^\infty(0,T;L^1(\mathbb{R}))$，在\eqref{align:ch07-7.161}中令 $j\to\infty$，即得\eqref{align:ch07-7.17}。
\end{proof}

\begin{proof}[\cref{proposition:ch07-7.2.1}的证明]
我们证明如下结论。

\paragraph{结论 1.}对几乎所有 $t\in(0,T)$，映射 $x\mapsto Y(x,t)\triangleq \int_0^t\phi(u_r(x,s))\,\mathrm{d}s$ 属于 $BV(\mathbb{R})$；且对任意 $a,b\in\mathbb{R}$，$a\leqslant b$，有
\begin{align}
u(\cdot,t)([a,b])-u_0([a,b])=Y(a^-,0,t)-Y(b^+,0,t). \label{align:ch07-7.162}
\end{align}

取 $a=b=x_0$。由\eqref{align:ch07-7.162}，对几乎所有 $0\leqslant t\leqslant T$，
\begin{align*}
u_s(\cdot,t)(\{x_0\})=u_{0s}(\{x_0\})+Y(x_0^-,0,t)-Y(x_0^+,0,t)
\geqslant u_{0s}(\{x_0\})-2\|\phi\|_\infty t,
\end{align*}
故当 $t\in(0,u_{0s}(\{x_0\})/(2\|\phi\|_\infty))$ 时 $u_s(\cdot,t)(\{x_0\})>0$。由此得\eqref{align:ch07-7.19}。

余下证结论。如\eqref{align:ch07-7.161}那样论证，得
\begin{align}
\int_\mathbb{R}Y(x,\tau)\rho'(x)\,\mathrm{d}x=\langle u(\cdot,\tau),\rho\rangle-\langle u_0,\rho\rangle. \label{align:ch07-7.163}
\end{align}
故分布导数 $\partial_xY(x,t)$ 属于 $\mathscr{R}_f(\mathbb{R})$。另一方面，由\eqref{align:ch07-7.1}，$|Y(\cdot,t)|\leqslant M\int_0^t u_r(\cdot,s)\,\mathrm{d}s\in L^1(\mathbb{R})$，故对几乎所有 $t\in(0,T)$，$Y(\cdot,t)\in BV(\mathbb{R})$。

为证\eqref{align:ch07-7.162}，由标准正则化论证可在\eqref{align:ch07-7.163}中取 $\rho=\rho_j$，
\begin{align*}
\rho_j(x)\triangleq j\left(x-a+\frac1j\right)\chi_{[a-1/j,a]}(x)+\chi_{(a,b)}(x)+j\left(b+\frac1j-x\right)\chi_{[b,b+1/j]}(x),
\end{align*}
其中 $a,b\in\mathbb{R}$，$a<b$，$j\in\mathbb{N}$。于是
\begin{align}
\langle u(\cdot,\tau),\rho_j\rangle-\langle u_0,\rho_j\rangle
=j\int_{a-1/j}^a Y(x,\tau)\,\mathrm{d}x-j\int_b^{b+1/j}Y(x,\tau)\,\mathrm{d}x. \label{align:ch07-7.164}
\end{align}
在\eqref{align:ch07-7.164}中令 $j\to\infty$，即得\eqref{align:ch07-7.162}。结论得证。
\end{proof}

\begin{remark}\label{remark:ch07-7.6.1}
\cref{proposition:ch07-7.2.1}的证明事实上依赖于一维双曲守恒律的熵解与 Hamilton--Jacobi 方程的粘性解之间的对应（参见文献~\cite{ref21,ref20}）。
\end{remark}

不妨在 (A2') 成立且 $\phi''<0$、$\phi'>0$ 在 $(0,\infty)$、$\phi(0)=0$（从而 $\alpha\phi+\beta>0$，$\beta>0$，且 (A1) 亦满足；见\cref{remark:ch07-7.1.1}）的假设下证明\cref{theorem:ch07-7.3.3}与\cref{proposition:ch07-7.3.5}。考虑 (Pn) 的下列正则化问题（与\ref{section:ch07-逼近问题}中所用的 (P$n^\epsilon$) 不同）：
\begin{align*}
\begin{cases}
\partial_t y_n^\epsilon+\partial_x[\phi(y_n^\epsilon)]=\epsilon\,\partial_{xx}[\phi(y_n^\epsilon)] & \text{在 }S,\\
y_n^\epsilon=u_{0n}^\epsilon & \text{在 }\mathbb{R}\times\{0\},
\end{cases}
\qquad(Vn^\epsilon)
\end{align*}
其中 $\{u_{0n}^\epsilon\}$ 满足\eqref{align:ch07-7.46}--\eqref{align:ch07-7.47}。\ref{section:ch07-逼近问题}中为 (P$n^\epsilon$) 回顾的存在唯一与正则性结果、\cref{lemma:ch07-7.4.1}的先验估计及\cref{lemma:ch07-7.4.5}的收敛结果，对 (V$n^\epsilon$) 的解同样成立（见文献~\cite{ref63}）。族 $\{y_n^\epsilon\}$ 的若干性质是下列两个引理的内容。

\begin{lemma}\label{lemma:ch07-7.6.1}
设 (A2') 成立，$\phi(0)=0$，$\phi''<0$、$\phi'>0$ 在 $(0,\infty)$。设 $u_n$ 为\cref{proposition:ch07-7.4.7}给出的问题 (Pn) 的唯一熵解。则存在 (V$n^\epsilon$) 解的子列 $\{y_n^{\epsilon_j}\}$，使
\begin{align}
y_n^{\epsilon_j}\stackrel{*}{\rightharpoonup}u_n\quad\text{在 }L^\infty(S),\qquad
y_n^{\epsilon_j}\to u_n\quad\text{在 }L^1(R_L)\ \text{对所有 }L>0. \label{align:ch07-7.165}
\end{align}
\end{lemma}

\begin{proof}
由\eqref{align:ch07-7.66}与\eqref{align:ch07-7.69}，存在序列 $\{y_n^{\epsilon_j}\}$ 与 $y_n\in L^\infty(S)\cap L^\infty(0,T;L^1(\mathbb{R}))$ 使\eqref{align:ch07-7.165}成立。此外，对任意凸函数 $E$、满足 $F'=E'\phi'$ 的函数 $F$ 及如\cref{definition:ch07-Pn熵解}的 $\zeta$，有
\begin{align*}
\iint_S\{E(y_n^{\epsilon_j})\partial_t\zeta+F(y_n^{\epsilon_j})\partial_x\zeta\}\,\mathrm{d}x\,\mathrm{d}t
+\int_\mathbb{R}E(u_{0n}^{\epsilon_j})\zeta(x,0)\,\mathrm{d}x
\geqslant\epsilon_j\iint_S F'(y_n^{\epsilon_j})\partial_{xy}y_n^{\epsilon_j}\partial_x\zeta\,\mathrm{d}x\,\mathrm{d}t.
\end{align*}
令 $j\to\infty$ 并如\cref{proposition:ch07-7.4.7}证明中那样论证，得
\begin{align*}
\iint_S[E(y_n)\partial_t\zeta+F(y_n)\partial_x\zeta]\,\mathrm{d}x\,\mathrm{d}t
\geqslant-\int_\mathbb{R}E(u_{0n})\zeta(x,0)\,\mathrm{d}x.
\end{align*}
故 $y_n$ 是 (Pn) 的熵解，由 Kružkov 唯一性定理得 $y_n=u_n$。结论得证。
\end{proof}

\begin{lemma}\label{lemma:ch07-7.6.2}
设 (A0) 与 (A2') 成立，$\phi(0)=0$，$\phi''<0$、$\phi'>0$ 在 $(0,\infty)$。则对所有 $t\in(0,T)$、$\epsilon>0$ 与 $n\in\mathbb{N}$，有
\begin{align}
\pm\frac{\partial}{\partial t}\left[\frac{\alpha\phi(y_n^\epsilon(\cdot,t))+\beta}{t^\alpha}\right]\leqslant0\quad\text{在 }\mathbb{R}\text{ 若 }\pm\alpha>0, \label{align:ch07-7.166a}
\end{align}
\begin{align}
\frac{\partial}{\partial t}[\phi(y_n^\epsilon(\cdot,t))-\beta\log t]\leqslant0\quad\text{在 }\mathbb{R}\text{ 若 }\alpha=0. \label{align:ch07-7.166b}
\end{align}
\end{lemma}

\begin{proof}
为方便置 $A\triangleq \epsilon\partial_{xx}-\partial_x$，于是 $\partial_t y_n^\epsilon=A[\phi(y_n^\epsilon)]$ 在 $S$。令
\begin{align*}
z_n^\epsilon\triangleq t\partial_t y_n^\epsilon-g(y_n^\epsilon),\qquad
g(y_n^\epsilon)\triangleq \frac{\alpha\phi(y_n^\epsilon)+\beta}{\phi'(y_n^\epsilon)}\quad(n\in\mathbb{N}).
\end{align*}
由 (A2') 与直接计算得
\begin{align*}
\partial_t z_n^\epsilon=A[\phi'(y_n^\epsilon)z_n^\epsilon]
+\frac{[\alpha+1-g'(y_n^\epsilon)][z_n^\epsilon+g(y_n^\epsilon)]}{t}
\leqslant A[\phi'(y_n^\epsilon)z_n^\epsilon]+\frac{[\alpha+1-g'(y_n^\epsilon)]z_n^\epsilon}{t}
\quad\text{在 }S.
\end{align*}
因 $z_n^\epsilon(\cdot,0)=-g(u_{0n}^\epsilon)\leqslant0$ 在 $\mathbb{R}$，由比较原理得 $z_n^\epsilon\leqslant0$ 在 $S$。从而 $t\partial_t y_n^\epsilon(\cdot,t)\leqslant g(y_n^\epsilon)(\cdot,t)$ 在 $\mathbb{R}$ 对所有 $t\in(0,T)$，由此得\eqref{align:ch07-7.166a}。
\end{proof}

为证\cref{theorem:ch07-7.3.3}(i)，还需要一个引理。

\begin{lemma}\label{lemma:ch07-7.6.3}
设 (A0)--(A1) 成立，$u$ 为所构造解。设 $\{u_{n_k}\}$ 为\cref{theorem:ch07-7.5.1}证明中所用的序列。则对几乎所有 $t\in(0,T)$ 与每个 $x_0\in\operatorname{supp}u_s(\cdot,t)$，存在序列 $\{x_p\}\subseteq\mathbb{R}$ 与子列 $\{u_{n_p}\}\subseteq\{u_{n_k}\}$，使 $x_p\to x_0$ 且 $u_{n_p}(x_p,t)\to\infty$（当 $p\to\infty$）。
\end{lemma}

\begin{proof}
由\cref{remark:ch07-7.5.1}，存在零测集 $N\subseteq(0,T)$ 与 $\{u_{n_k}\}$ 的子列（不重新标号），使
\begin{align}
u_{n_k}(\cdot,t)\stackrel{*}{\rightharpoonup}u_r(\cdot,t)+u_s(\cdot,t)\quad\text{在 }\mathscr{R}_f(I_L) \label{align:ch07-7.167}
\end{align}
对所有 $t\in(0,T)\setminus N$ 与 $L>0$ 成立（见\eqref{align:ch07-7.107}取 $U(y)=y$ 及\eqref{align:ch07-7.109a}--\eqref{align:ch07-7.109b}）。

设 $x_0\in\operatorname{supp}u_s(\cdot,t)$。则不存在邻域 $I_\delta(x_0)$ 使 $\{u_{n_k}(\cdot,t)\}$ 包含于 $L^\infty(I_\delta(x_0))$ 的有界子集；否则存在 $\{u_{n_k}(\cdot,t)\}$ 的子列（不重新标号）与 $f_t\in L^\infty(I_\delta(x_0))$、$f_t\geqslant0$，使 $u_{n_k}(\cdot,t)\stackrel{*}{\rightharpoonup}f_t$ 在 $L^\infty(I_\delta(x_0))$。然而由\eqref{align:ch07-7.167}，这将蕴含 $u_s(\cdot,t)=0$ 在 $I_\delta(x_0)$，矛盾。

由此对所有 $p\in\mathbb{N}$，$\sup_{n_k}\|u_{n_k}(\cdot,t)\|_{L^\infty(I_{1/p}(x_0))}=\infty$。于是存在 $\{u_{n_{k_p}}\}\subseteq\{u_{n_k}\}$ 与 $\{x_p\}$，使 $x_p\in I_{1/p}(x_0)$ 且 $u_{n_p}(x_p,t)\geqslant p$ 对所有 $p\in\mathbb{N}$。结论得证。
\end{proof}

\begin{proof}[\cref{theorem:ch07-7.3.3}的证明]
如前所述，只就 $\phi''<0$、$\phi'>0$ 在 $(0,\infty)$ 的情形证明。

(i) 设 $\phi$ 有界，$\{u_{n_k}\}$ 如\cref{lemma:ch07-7.6.3}证明中。由\cref{lemma:ch07-7.6.1}，对每个 $n_k$，存在 (V$n_{n_k}^\epsilon$) 解的子列 $\{y_{n_k}^{\epsilon_j}\}$，使对几乎所有 $t\in(0,T)$，
\begin{align}
y_{n_k}^{\epsilon_j}(\cdot,t)\to u_{n_k}(\cdot,t)\quad\text{在 }L^1_{\mathrm{loc}}(\mathbb{R})\ \text{当 }j\to\infty. \label{align:ch07-7.168}
\end{align}
另一方面，由\cref{lemma:ch07-7.6.2}的证明，对任意 $t\in(0,T)$，有
\begin{align}
\epsilon_j\partial_{xx}\phi(y_{n_k}^{\epsilon_j}(\cdot,t))-\partial_x\phi(y_{n_k}^{\epsilon_j}(\cdot,t))
=\partial_t y_{n_k}^{\epsilon_j}(\cdot,t)\leqslant\frac{g(y_{n_k}^{\epsilon_j}(\cdot,t))}{t}\quad\text{在 }\mathbb{R}, \label{align:ch07-7.169}
\end{align}
其中 $g(y)\triangleq (\alpha\phi(y)+\beta)/\phi'(y)>0$。

固定 $a,b\in\mathbb{R}$，$a<b$，设 $\rho\in C^1_c(a,b)$，$\rho\geqslant0$。置
\begin{align*}
Z(y)\triangleq \int_0^y\frac{[\phi'(z)]^2}{\alpha\phi(z)+\beta}\,\mathrm{d}z.
\end{align*}
则 $Z$ 递增，且 $0\leqslant Z(y)\leqslant\phi'(0)$ 对所有 $y\geqslant0$。将\eqref{align:ch07-7.169}乘以 $\rho/g(y_{n_k}^{\epsilon_j}(\cdot,t))$ 并积分，得
\begin{align*}
\int_a^b Z(y_{n_k}^{\epsilon_j})(x,t)[\epsilon_j\rho''(x)+\rho'(x)]\,\mathrm{d}x
\leqslant\frac1t\int_a^b\rho(x)\,\mathrm{d}x
-\epsilon_j\int_a^b\frac{\phi'(y_{n_k}^{\epsilon_j})g'(y_{n_k}^{\epsilon_j})}{[g(y_{n_k}^{\epsilon_j})]^2}[\partial_x y_{n_k}^{\epsilon_j}]^2(x,t)\rho(x)\,\mathrm{d}x
\leqslant\frac1t\int_a^b\rho(x)\,\mathrm{d}x.
\end{align*}
由\eqref{align:ch07-7.168}，令 $j\to\infty$，得
\begin{align}
\int_a^b Z(u_{n_k}(x,t))\rho'(x)\,\mathrm{d}x\leqslant\frac1t\int_a^b\rho(x)\,\mathrm{d}x. \label{align:ch07-7.170}
\end{align}
对几乎所有 $t\in(0,T)$，设 $x_0\in\operatorname{supp}u_s(\cdot,t)$。证明
\begin{align}
\operatorname*{ess\,lim}_{x\to x_0^+}u_r(x,t)=\infty. \label{align:ch07-7.171}
\end{align}
设 $\{x_p\}$ 与 $\{u_{n_p}\}\subseteq\{u_{n_k}\}$ 由\cref{lemma:ch07-7.6.3}给出，固定 $b>x_0$。因 $x_p\to x_0$（当 $p\to\infty$），存在 $\bar p\in\mathbb{N}$ 使 $b>x_p$ 对所有 $p>\bar p$。对任意此种 $p$ 与充分大的 $m\in\mathbb{N}$，置
\begin{align*}
\rho_{p,m}(x)\triangleq m(x-x_p)\chi_{[x_p,x_p+1/m]}(x)+\chi_{(x_p+1/m,b-1/m)}(x)+m(b-x)\chi_{[b-1/m,b]}(x).
\end{align*}
在\eqref{align:ch07-7.170}中取 $a=x_p$、$\rho=\rho_{p,m}$，得
\begin{align*}
m\int_{x_p}^{x_p+1/m}Z(u_{n_p}(x,t))\,\mathrm{d}x-m\int_{b-1/m}^b Z(u_{n_p}(x,t))\,\mathrm{d}x
\leqslant\frac1t\int_{x_p}^b\rho_{p,m}(x)\,\mathrm{d}x,
\end{align*}
从而令 $m\to\infty$，对所有 $p>\bar p$ 有
\begin{align*}
Z(u_{n_p}(x_p,t))\leqslant Z(u_{n_p}(b,t))+\frac{b-x_p}{t}.
\end{align*}
由\cref{lemma:ch07-7.6.3}证明，$u_{n_p}(x_p,t)\geqslant p$。令 $p\to\infty$ 并利用\eqref{align:ch07-7.143}，对几乎所有 $b>x_0$，得
\begin{align}
\lim_{y\to\infty}Z(y)\leqslant Z(u_r(b,t))+\frac{b-x_0}{t}, \label{align:ch07-7.172}
\end{align}
从而
\begin{align*}
\lim_{y\to\infty}Z(y)\leqslant\operatorname*{ess\,lim}_{b\to x_0^+}Z(u_r(b,t)).
\end{align*}
因 $Z$ 连续且递增，由上式得\eqref{align:ch07-7.171}。

现可证 $u$ 在 $[0,\tau]$ 于 $x_0\in\operatorname{supp}u_s(\cdot,\tau)$ 处满足相容条件。关于 (C$^-$)，由 $\phi$ 的递增性：事实上若 $u_r>l$ 则 $H_-(u_r-l)=0$；若 $u_r\leqslant l$，则
\begin{align*}
H_-(u_r(x,t)-l)[\phi(u_r(x,t))-\phi(l)]=\phi(l)-\phi(u_r(x,t))\geqslant0,
\end{align*}
故 (C$^-$) 满足。关于 (C$^+$)，由等式\eqref{align:ch07-7.171}，对所有 $l\in[0,\infty)$，对几乎所有 $t\in(0,\tau)$ 有 $\operatorname*{ess\,lim}_{x\to x_0^+}H_-(u_r(\cdot,t)-l)=0$（回忆由单调性此类 $t$ 有 $x_0\in\operatorname{supp}u_s(\cdot,t)$）。于是由控制收敛定理得 (C$^+$)。

(ii) 设 $\phi$ 无界。由\eqref{align:ch07-7.166a}，对所有 $0<t_1\leqslant t\leqslant T$ 与 $x\in\mathbb{R}$，
\begin{align*}
\int_{t_1}^t\phi(y_{n_k}^{\epsilon_j})(x,s)\,\mathrm{d}s
\geqslant
\begin{cases}
\dfrac{\alpha\phi(y_{n_k}^{\epsilon_j})(x,t)+\beta}{\alpha t^\alpha}\dfrac{t^{\alpha+1}-t_1^{\alpha+1}}{\alpha+1}-\dfrac{\beta}{\alpha}(t-t_1)&\text{若 }\alpha\neq0,\\[2mm]
[\phi(y_{n_k}^{\epsilon_j})(x,t)-\beta\log t](t-t_1)+\beta[t\log t-t]-\beta[t_1\log t_1-t_1]&\text{若 }\alpha=0,
\end{cases}
\end{align*}
其中 $\{y_{n_k}^{\epsilon_j}\}$ 为\cref{lemma:ch07-7.6.1}证明中所用子列。令 $\epsilon_j\to0$，由\eqref{align:ch07-7.165}，对几乎所有 $t\in(t_1,T)$ 与 $x\in\mathbb{R}$ 得
\begin{align}
\int_{t_1}^t\phi(u_{n_k}(x,s))\,\mathrm{d}s\geqslant
\begin{cases}
\dfrac{\alpha\phi(u_{n_k})(x,t)+\beta}{\alpha t^\alpha}\dfrac{t^{\alpha+1}-t_1^{\alpha+1}}{\alpha+1}-\dfrac{\beta}{\alpha}(t-t_1)&\text{若 }\alpha\neq0,\\
[\phi(u_{n_k})(x,t)-\beta](t-t_1)+\beta t_1\log\dfrac{t}{t_1}&\text{若 }\alpha=0.
\end{cases} \label{align:ch07-7.173}
\end{align}
另一方面，如\cref{proposition:ch07-7.2.1}证明中那样，由弱形式\eqref{align:ch07-7.78}易见，对每个 $t\in(0,T]$，映射 $x\mapsto Y_n(x)\triangleq \int_0^t\phi(u_{n_k}(x,s))\,\mathrm{d}s$ 属于 $W^{1,1}(\mathbb{R})$，其弱导数 $Y_n'(\cdot)=u_{0n}-u_{n_k}(\cdot,t)$ 在 $\mathbb{R}$ 中几乎处处。故
\begin{align*}
0\leqslant Y_n\leqslant\|u_{0n}\|_{L^1(\mathbb{R})}\leqslant\|u_0\|_{\mathscr{R}_f(\mathbb{R})}.
\end{align*}
结合此不等式与\eqref{align:ch07-7.173}并令 $t_1\to0^+$，得
\begin{align}
\phi(u_{n_k})(x,t)\leqslant\frac{(\alpha+1)\|u_0\|_{\mathscr{R}_f(\mathbb{R})}}{t}+\beta
\quad\text{对几乎所有 }t\in(0,T)\text{ 与 }x\in\mathbb{R} \label{align:ch07-7.174}
\end{align}
（回忆对无界 $\phi$，因假设 $\phi'>0$ 且 $\alpha\phi+\beta>0$ 在 $[0,\infty)$，故 $\alpha\geqslant0$）。由\eqref{align:ch07-7.174}，序列 $\{u_{n_k}\}$ 对每个 $\tau\in(0,T)$ 包含于 $L^\infty(\mathbb{R}\times(\tau,T))$ 的有界子集，于是由\cref{theorem:ch07-7.3.1}，令 $n\to\infty$ 即得结论。
\end{proof}

\begin{proof}[\cref{proposition:ch07-7.3.5}的证明]
由\cref{theorem:ch07-7.3.3}的证明，对几乎所有 $t\in(0,T)$，不等式\eqref{align:ch07-7.172}对所有 $x_j\in\operatorname{supp}u_s(\cdot,t)$ 成立。设 $x_1\in\operatorname{supp}u_s(\cdot,t)$，置 $I_1\triangleq (x_1-\epsilon,x_1+\epsilon)$（$\epsilon>0$）。因 $Z$ 可逆，由\eqref{align:ch07-7.172}得
\begin{align*}
\int_{I_1}u_r(x,t)\,\mathrm{d}x\geqslant\int_{x_1}^{x_1+\epsilon}Z^{-1}\!\left(\lim_{y\to\infty}Z(y)-\frac{x-x_1}{t}\right)\mathrm{d}x
=\int_0^\epsilon Z^{-1}\!\left(\lim_{y\to\infty}Z(y)-\frac{y}{t}\right)\mathrm{d}y=:B_\epsilon,
\end{align*}
其中
\begin{align}
\lim_{\epsilon\to0^+}\frac{B_\epsilon}{\epsilon}=Z^{-1}\!\left(\lim_{y\to\infty}Z(y)\right)=\infty. \label{align:ch07-7.175}
\end{align}
若 $\operatorname{supp}u_s(\cdot,t)\not\subseteq I_1$，取 $x_2\in\operatorname{supp}u_s(\cdot,t)\setminus I_1$，置 $I_2\triangleq (x_2-\epsilon,x_2+\epsilon)$。因 $(x_1,x_1+\epsilon)\cap(x_2,x_2+\epsilon)=\varnothing $，
\begin{align*}
\int_{I_1\cup I_2}u_r(x,t)\,\mathrm{d}x\geqslant\int_{x_1}^{x_1+\epsilon}u_r\,\mathrm{d}x+\int_{x_2}^{x_2+\epsilon}u_r\,\mathrm{d}x\geqslant2B_\epsilon.
\end{align*}
递归地，只要 $\operatorname{supp}u_s(\cdot,t)\not\subseteq\bigcup_{i=1}^{n-1}I_i$（$I_i\triangleq (x_i-\epsilon,x_i+\epsilon)$），就存在 $x_n\in\operatorname{supp}u_s(\cdot,t)\setminus\bigcup_{i=1}^{n-1}I_i$，使
\begin{align*}
nB_\epsilon\leqslant\int_{\bigcup_{i=1}^n I_i}u_r(x,t)\,\mathrm{d}x\leqslant\|u_0\|_{\mathscr{R}_f(\mathbb{R})}.
\end{align*}
故过程在某个 $n=n_\epsilon$ 处终止，$n_\epsilon B_\epsilon\leqslant\|u_0\|_{\mathscr{R}_f(\mathbb{R})}$。于是
\begin{align*}
\operatorname{supp}u_s(\cdot,t)\subseteq\bigcup_{i=1}^{n_\epsilon}I_i,
\end{align*}
从而
\begin{align*}
\lambda(\operatorname{supp}u_s(\cdot,t))\leqslant\lambda\!\left(\bigcup_{i=1}^{n_\epsilon}I_i\right)\leqslant2n_\epsilon\epsilon
\leqslant\frac{2\|u_0\|_{\mathscr{R}_f(\mathbb{R})}\epsilon}{B_\epsilon}.
\end{align*}
令 $\epsilon\to0^+$ 并利用\eqref{align:ch07-7.175}，即得结论。
\end{proof}




\end{document}